\chapter{双编码器系统工作原理}
\label{ch:dual_encoder}

\section{双编码器的真实价值}

大多数控制工程师对双编码器的理解停留在："一个电机端测速、一个输出端定位"。\textbf{这是完全错误的认知。}

\begin{keyinsight}{双编码器的真正价值}
双编码器的核心价值是：\\
\textbf{通过两端角度差计算关节的扭转形变，从而间接得到关节力矩。}\\
它不是为了提高定位精度，而是为了\textbf{辅助力矩观测}。
\end{keyinsight}

\section{双编码器配置方案}

\subsection{电机端编码器}

\begin{itemize}
  \item \textbf{位置}：安装在电机轴（减速器输入侧）
  \item \textbf{功能}：
  \begin{itemize}
    \item 电机转子位置$\theta_m$——用于\textbf{FOC矢量控制}的Park/Clarke变换
    \item 电机转速$\omega_m$——用于速度环控制
    \item 记录磁极位置——用于\textbf{电机换向}
  \end{itemize}
  \item \textbf{典型规格}：增量式/绝对值，分辨率$2^{14}$--$2^{19}$脉冲/转
  \item \textbf{核心要求}：高更新率（通常$>10$ kHz）、低延迟
\end{itemize}

\subsection{输出端编码器}

\begin{itemize}
  \item \textbf{位置}：安装在关节输出轴（减速器输出侧）
  \item \textbf{功能}：
  \begin{itemize}
    \item 关节绝对位置$\theta_j$——用于运动学计算和位置控制
    \item 与电机端角度对比$\Delta\theta = \theta_m/i - \theta_j$——计算\textbf{传动系统形变}
    \item 关节长期定位精度——补偿减速器的累积误差
  \end{itemize}
  \item \textbf{典型规格}：绝对值编码器，分辨率$2^{17}$--$2^{23}$脉冲/转
  \item \textbf{核心要求}：低温漂、抗振动
\end{itemize}

\begin{figure}[H]
\centering
\begin{tcolorbox}[colback=lightgray]
\centering
\textbf{双编码器系统信号流}\\
\vspace{5pt}
$\theta_m$（电机端）$\underset{\div i}{\longrightarrow}$ $\theta_m/i$（电机端折算到输出）\\
$\theta_j$（输出端）\\
$\Delta\theta = \theta_m/i - \theta_j$（扭转角）$\underset{\times K}{\longrightarrow}$ 关节力矩估计
\end{tcolorbox}
\caption{双编码器力矩估计原理}
\label{fig:dual_encoder_principle}
\end{figure}

\section{通过双编码器计算关节力矩}

\subsection{基本原理}

减速器和关节结构可以简化为一个\textbf{弹性体}（扭转弹簧）：

\begin{equation}
T_{\text{joint}} = K_{\text{stiffness}} \times (\theta_m/i - \theta_j) = K_{\text{stiffness}} \times \Delta\theta
\end{equation}

其中$K_{\text{stiffness}}$是关节传动系统（减速器+结构件）的等效扭转刚度。

\begin{engineeringbox}{双编码器力矩估计的精度}
力矩估计精度取决于：\\
\begin{itemize}
  \item 编码器分辨率（决定$\Delta\theta$的测量精度）
  \item 刚度$K_{\text{stiffness}}$的准确性和线性度
  \item 温度对刚度的影响（柔性材料$K$随温度变化）
\end{itemize}
在良好的校准下，双编码器力矩估计精度可达\textbf{2--5\%额定力矩}，\\接近低精度力矩传感器的水平。
\end{engineeringbox}

\subsection{刚度$K_{\text{stiffness}}$的标定方法}

$K_{\text{stiffness}}$是双编码器力矩估计中\textbf{最关键也最容易引入误差}的参数。以下是获取$K_{\text{stiffness}}$的几种途径。

\subsubsection{方法一：从减速器手册查表}

这是最直接的方式。减速器厂商通常会在规格书中给出扭转刚度：

\begin{table}[H]
\centering
\begin{tabular}{lc}
\toprule
\textbf{减速器类型} & \textbf{典型刚度范围（N·m/rad）} \\
\midrule
RV减速器 & $10^5 \sim 10^6$ \\
谐波减速器 & $10^3 \sim 10^5$ \\
行星减速器 & $10^4 \sim 10^5$ \\
\bottomrule
\end{tabular}
\caption{常见减速器扭转刚度典型值}
\label{tab:typical_stiffness}
\end{table}

以谐波减速器CSF-25为例，其扭转刚度约为$1.2 \times 10^4$ N·m/rad；\newline
RV-40E约为$2.5 \times 10^5$ N·m/rad。实际值需查阅对应型号的手册。

\subsubsection{方法二：实验标定}

获取$K_{\text{stiffness}}$最准确的方式是实验标定：

\begin{enumerate}
  \item 锁定关节输出端（使$\theta_j = 0$）
  \item 在电机端施加一系列已知力矩$T_m$
  \item 读取双编码器，记录每次的$\Delta\theta = \theta_m/i - \theta_j$
  \item 绘制$T_m$对$\Delta\theta$的散点图，进行线性拟合
\end{enumerate}

\begin{equation}
T_m = K_{\text{stiffness}} \cdot \Delta\theta + T_{\text{friction}}
\label{eq:stiffness_fit}
\end{equation}

拟合得到的\textbf{斜率}即为$K_{\text{stiffness}}$。

\begin{cautionbox}{摩擦力补偿}
低速时摩擦会吃掉一部分力矩，使得单次测量偏差较大。\\
建议在多个不同电流值（力矩）下采集数据点，\\通过\textbf{最小二乘线性拟合}消除摩擦力的影响，\\而不是仅靠单点测量。
\end{cautionbox}

\subsubsection{方法三：归一化仿真假设}

在研发早期或仿真阶段，如果无法获取精确的$K_{\text{stiffness}}$，\newline
可以先采用一个合理假设值观察柔性效应，待原型机完成后重新标定。

\section{电机力矩常数$K_t$的获取}

在标定$K_{\text{stiffness}}$的实验步骤中，\textbf{"施加已知力矩$T_m$"}指的是电机的电磁力矩。\newline
它通过电机力矩常数$K_t$和q轴电流$I_q$计算得到：

\begin{equation}
T_m = K_t \times I_q
\label{eq:torque_current}
\end{equation}

因此，$K_t$的准确性直接决定了$K_{\text{stiffness}}$的标定精度。

\subsection{查手册获取}

电机制造商通常会在规格书中给出力矩常数，单位一般为N·m/A或mN·m/A：

\begin{equation}
K_t = 0.12\ \text{N·m/A} \quad \text{（示例值）}
\label{eq:kt_example}
\end{equation}

\subsection{从额定参数反算}

如果手册中没有直接给出$K_t$，可以通过额定参数估算：

\begin{equation}
K_t = \frac{60 \times P_{\text{rated}}}{2\pi \times n_{\text{rated}} \times I_{\text{rated}}}
\label{eq:kt_calc}
\end{equation}

其中：
\begin{itemize}
  \item $P_{\text{rated}}$——额定功率（W）
  \item $n_{\text{rated}}$——额定转速（rpm）
  \item $I_{\text{rated}}$——额定电流（A）
  \item $K_t$——力矩常数（N·m/A）
\end{itemize}

\subsection{关于转矩常数与反电动势常数的关系}

对于永磁同步电机，在SI单位制下，力矩常数$K_t$与反电动势常数$K_e$在数值上相等：

\begin{equation}
K_t\ [\text{N·m/A}] = K_e\ [\text{V·s/rad}]
\label{eq:kt_ke}
\end{equation}

因此，如果手册给出了反电动势常数，也可以直接作为$K_t$使用。

\subsection{动力学滤波}

直接差分得到的$\Delta\theta$含有噪声，需要滤波：

\begin{equation}
\hat{T}_{\text{joint}}(s) = \frac{K_{\text{stiffness}} \cdot \Delta\theta(s)}{\tau s + 1}
\end{equation}

其中$\tau$是低通滤波时间常数。滤波会引入相位滞后，在动态任务中需要补偿。

\section{双编码器如何补偿减速器误差}

\subsection{减速器的主要可补偿误差}

双编码器可以补偿的误差类型：

\begin{enumerate}
  \item \textbf{位置依赖的传动误差（PDE）}：
  \begin{itemize}
    \item 减速器内部齿轮啮合不均匀导致的位置误差
    \item 可以通过$\Delta\theta(\theta_j)$映射表补偿
  \end{itemize}
  \item \textbf{滞后}：
  \begin{itemize}
    \item 与力矩方向相关的传动误差
    \item 双编码器可以检测到"电机转了但输出没动"的背隙区域
  \end{itemize}
  \item \textbf{温度漂移}：
  \begin{itemize}
    \item 热膨胀导致减速器齿隙变化
    \item 双编码器实时检测，不需要温度模型
  \end{itemize}
\end{enumerate}

\subsection{双编码器不能补偿的误差}

\begin{itemize}
  \item \textbf{摩擦}：摩擦是力矩损失，不是位置偏差，双编码器检测不到
  \item \textbf{高频力矩脉动}：高于编码器采样率的脉动，无法被差分检测
  \item \textbf{结构共振}：$\Delta\theta$检测的是静态/准静态形变，不是动态振动
\end{itemize}

\section{半直驱关节不需要力矩传感器的原理}

这是本书\textbf{另一个核心问题}的答案：

半直驱关节（$i \approx 6\text{--}15$）之所以\textbf{可以不用力矩传感器}，原因有二：

\subsection{原因一：电流力矩观测可用}

如第4章所述，低减速比下电流力矩观测精度5--10\%，满足力控需求。不需要额外的力矩传感器。

\subsection{原因二：双编码器提供交叉验证}

\begin{itemize}
  \item 通过电流估算的力矩 $T_{\text{current}} = K_t \cdot i_q \cdot i$
  \item 通过双编码器估算的力矩 $T_{\text{dual}} = K_{\text{stiffness}} \cdot \Delta\theta$
  \item 两个独立估计可以互相验证和融合
  \item 当两者偏差超出阈值时，可触发故障检测
\end{itemize}

\begin{equation}
T_{\text{final}} = \alpha \cdot T_{\text{current}} + (1-\alpha) \cdot T_{\text{dual}}
\end{equation}

其中$\alpha$是融合权重，根据工况动态调整。

\begin{table}[H]
\centering
\begin{tabular}{lccc}
\toprule
\textbf{方案} & \textbf{电流力控} & \textbf{双编码器} & \textbf{力矩传感器} \\
\midrule
宇树H1 & ✓ 低减速比可行 & ✓ 辅助 & ✗ 不需要 \\
特斯拉Optimus & ✗ 高减速比失效 & ✓ 必须 & ✓ 必须 \\
优必选Walker & ✗ 超高减速比 & ✓ 必须 & ✓ 必须 \\
传统工业臂 & ✗ & ✗ 只有电机端编码器 & ✓ 末端 \\
\bottomrule
\end{tabular}
\caption{主流关节方案传感器配置}
\label{tab:sensor_config}
\end{table}

\section{双编码器的选型注意事项}

\subsection{分辨率匹配}

\begin{equation}
\Delta\theta_{\text{min}} = \frac{T_{\text{resolution}}}{K_{\text{stiffness}}}
\end{equation}

如果需要的力矩分辨率是$T_{\text{resolution}}$，则编码器需要能检测到的最小$\Delta\theta$是上式。

\textbf{例子}：如果$K_{\text{stiffness}} = 5000$ Nm/rad，需要0.1 Nm的力矩分辨率：
\begin{equation}
\Delta\theta_{\text{min}} = \frac{0.1}{5000} = 2 \times 10^{-5} \text{ rad} \approx 0.0011^\circ
\end{equation}

需要输出端编码器分辨率至少为$2^{19}$脉冲/转才能可靠检测。

\subsection{更新率匹配}

\begin{itemize}
  \item 电机端编码器更新率必须匹配\textbf{电流环频率}（通常10--40 kHz）
  \item 输出端编码器更新率匹配\textbf{力控频率}（通常1--5 kHz）
  \item 两个编码器的采样\textbf{必须同步}——否则$\Delta\theta$的噪声被采样时间差淹没
\end{itemize}

\begin{cautionbox}{编码器不同步的陷阱}
如果电机端和输出端编码器不是同时采样，\\
即使实际的$\Delta\theta$为0，由于采样时间差和运动速度，\\
计算出的$\Delta\theta$也会非零：$\Delta\theta_{\text{false}} = \omega \times \Delta t_{\text{sampling}}$\\
这个\textbf{假的扭转角}会直接变成力矩估计的\textbf{虚假信号}。
\end{cautionbox}

\section{三大机器人路线的对比总结}

\begin{table}[H]
\centering
\begin{tabular}{p{2.5cm}p{3cm}p{3cm}p{3cm}}
\toprule
\textbf{特性} & \textbf{宇树路线} & \textbf{特斯拉路线} & \textbf{传统路线} \\
\midrule
减速比 & 低（6--15） & 高（$\approx$100） & 高（50--160） \\
电机端编码器 & 有 & 有 & 有 \\
输出端编码器 & 有 & 有 & \textbf{无} \\
力矩传感器 & \textbf{无} & \textbf{有}（关节端） & 有（腕部/末端） \\
力控实现 & 电流+双编码器 & 力矩传感器 & 力矩传感器 \\
力控精度 & 中（5--10\%） & 高（$<1\%$） & 高 \\
成本 & 低 & 高 & 高 \\
\bottomrule
\end{tabular}
\caption{三大机器人路线传感器与力控方案}
\label{tab:three_approaches}
\end{table}

\section*{本章小结}
\addcontentsline{toc}{section}{本章小结}

\begin{summarybox}{}
\begin{enumerate}
  \item 双编码器的核心价值是\textbf{通过$\Delta\theta$计算关节力矩}，不是提高定位精度
  \item 电机端编码器$\rightarrow$FOC换向，输出端编码器$\rightarrow$关节位置+力矩估计
  \item 半直驱关节（低减速比）\textbf{可以不用力矩传感器}——电流观测+双编码器足够
  \item 高减速比关节\textbf{必须用力矩传感器}——电流观测已失效，双编码器不够准
  \item 双编码器力矩估计精度取决于：分辨率、刚度标定、采样同步
  \item 三条技术路线的最核心差异在于：\textbf{力控信息靠什么获取}
\end{enumerate}
\end{summarybox}

\newpage
